% !PRINTED 342
\weberpage{342}
%<<CONT
mit ganzzahligen $t_1, t_2$ darstellbar sind (Bd.\,II, §\,163, 164). Denn
dann ist, wenn $\alpha$ durch $\lambda$ darstellbar ist, auch jedes Produkt $\omega\alpha$
durch $\lambda$ darstellbar. Für unseren Fall reduziert sich diese Bedingung darauf, daß $\alpha_1\theta, \alpha_2\theta$ durch $\lambda$ darstellbar sein müssen. Ist
aber $(a, b, c)$ irgend eine quadratische Form der Diskriminanten $\Delta$
mit \gesperrt{positivem ersten} Koeffizienten $a$, so können wir für $\theta$ auch
$\tfrac12(b+\sqrt{\Delta})$ nehmen und da
\[
\left(\frac{b+\sqrt{\Delta}}{2}\right)^2 = -ac + b\,\frac{b+\sqrt{\Delta}}{2}
\]
ist, so ist (11) immer die Basis eines Ideals in $\Omega$. Dabei kann
die positive Zahl $a_{22}$ willkürlich angenommen werden. Setzen wir
$a_{22} = 1$, so ist
\weq{15}{\lambda = a(x-\omega y)}
Basisform eines durch $(a, b, c)$ völlig bestimmten Ideals $\mathfrak{a}$, dessen
Norm gleich $a$ ist, und eine Basis dieses Ideals ist
\weq{16}{(a,\ a\omega).}

\begin{center}
§ 92. \gesperrt{Primideale im quadratischen Körper.}
\end{center}

Eine Primzahl $p$ ist im quadratischen Körper $\Omega$ entweder
selbst noch unzerlegbar und ist dann ein Primideal zweiten
Grades oder es zerfällt $p$ in zwei Primideale ersten Grades $\mathfrak{p}, \mathfrak{p}'$.
Wir haben hiernach zwei Fälle zu unterscheiden:
\[
\text{a)}\quad p = \mathfrak{p} \text{ in } \Omega \text{ unzerlegbar:}\quad N(\mathfrak{p}) = p^2,
\]
\[
\text{b)}\quad p = \mathfrak{p}\mathfrak{p}' \qquad\qquad N(\mathfrak{p}) = N(\mathfrak{p}') = p,
\]
und im letzten Falle können die beiden Primideale $\mathfrak{p}, \mathfrak{p}'$ entweder
voneinander verschieden oder auch gleich sein, so daß wir noch
einen dritten Fall unterscheiden können:
\[
\text{c)}\quad p = \mathfrak{p}^2,\quad N(\mathfrak{p}) = p.
\]
Es handelt sich nun darum, die Bedingungen zu ermitteln,
unter denen der eine oder der andere dieser Fälle eintritt.

Wenn $N(\mathfrak{p}) = p$ ist, so ist $p$ die Anzahl der inkongruenten
Zahlen $(\mathrm{mod}\ \mathfrak{p})$ in $\Omega$ und folglich ist in diesem Fall \gesperrt{jede Zahl}
in $\Omega$ einer der rationalen Zahlen $r$:
\weq{1}{0,\ 1,\ 2,\ \ldots,\ p-1}
kongruent.

Ist dagegen $N(\mathfrak{p}) = p^2$, so gibt es auch Zahlen, die nach $\mathfrak{p}$
nicht mit einer rationalen Zahl kongruent sind.
